\documentclass[11pt,letterpaper]{article} % lint-ok

% ── Geometry ──────────────────────────────────────────────────────────────────
\usepackage[margin=1in]{geometry}

% ── Pagination ────────────────────────────────────────────────────────────────
\usepackage[all]{nowidow}    % Require ≥2 lines on each side of a page break
\brokenpenalty=10000         % No hyphenated words across page breaks
\raggedbottom                % Allow uneven page bottoms rather than stretching

% ── Fonts ─────────────────────────────────────────────────────────────────────
\usepackage[T1]{fontenc}
\usepackage{newpxtext}     % Palatino-based serif (TeX Gyre Pagella)
\usepackage{newpxmath}     % Matching math font

% ── Colors ────────────────────────────────────────────────────────────────────
\usepackage[dvipsnames]{xcolor} % lint-ok
\definecolor{SectionBlue}{HTML}{1B3A5C}
\definecolor{AccentGray}{HTML}{4A4A4A}
\definecolor{QuoteBorder}{HTML}{8B9DAF}
\definecolor{RiskRed}{HTML}{C0392B}
\definecolor{RiskAmber}{HTML}{E67E22}
\definecolor{RiskGreen}{HTML}{27AE60}
\definecolor{BoxBg}{HTML}{F7F8FA}

% ── Tables ────────────────────────────────────────────────────────────────────
\usepackage{booktabs}
\usepackage{tabularx}

% ── Boxes ─────────────────────────────────────────────────────────────────────
\usepackage[framemethod=tikz]{mdframed}

% ── Lists ─────────────────────────────────────────────────────────────────────
\usepackage{enumitem}
\usepackage{needspace}

% ── Hyperlinks ────────────────────────────────────────────────────────────────
\usepackage{hyperref}
\usepackage{xurl}  % break bibliography URLs at any character, not just at slashes
\hypersetup{
  colorlinks=true,
  linkcolor=SectionBlue,
  urlcolor=SectionBlue,
  citecolor=SectionBlue,
  pdfauthor={JM Freeman},
  pdftitle={vox podcasts and audiobooks: PR/FAQ},
  pdfsubject={Working Backwards: Product Discovery},
}

% ── Bibliography ─────────────────────────────────────────────────────────────
\usepackage[style=numeric,backend=biber,sorting=none]{biblatex}
\addbibresource{prfaq-audio-programs.bib}

% ── Section styling ───────────────────────────────────────────────────────────
\usepackage{titlesec}
\titleformat{\section}
  {\needspace{6\baselineskip}\Large\bfseries\color{SectionBlue}}
  {}
  {0pt}
  {}

\titleformat{\subsection}
  {\needspace{5\baselineskip}\large\bfseries\color{AccentGray}}
  {}
  {0pt}
  {}

\titleformat{\subsubsection}
  {\normalsize\bfseries\color{AccentGray}}
  {}
  {0pt}
  {}
\titlespacing*{\subsubsection}{0pt}{1em}{0.5em}[0pt]

% ── Document stage ───────────────────────────────────────────────────────────
\newcommand{\prfaqstage}[1]{\def\theprfaqstage{#1}}
\prfaqstage{hypothesis}

% ── Document version ─────────────────────────────────────────────────────────
\newcommand{\prfaqversion}[2]{\def\theprfaqmajor{#1}\def\theprfaqminor{#2}}
\prfaqversion{1}{0}

% ── Header/footer ────────────────────────────────────────────────────────────
\usepackage{fancyhdr}
\pagestyle{fancy}
\fancyhf{}
\fancyhead[L]{\footnotesize\textcolor{AccentGray}{Stage: \theprfaqstage{} \enspace\textbar\enspace v\theprfaqmajor.\theprfaqminor}}
\fancyhead[R]{\footnotesize\textcolor{AccentGray}{\thepage}}
\renewcommand{\headrulewidth}{0pt}

% ══════════════════════════════════════════════════════════════════════════════
% Custom environments
% ══════════════════════════════════════════════════════════════════════════════

\newcommand{\prsection}[1]{%
  \vspace{1em}%
  {\large\bfseries\color{SectionBlue}#1}%
  \vspace{0.3em}\par%
}

\newmdenv[
  topline=false,
  bottomline=false,
  rightline=false,
  linewidth=3pt,
  linecolor=QuoteBorder,
  backgroundcolor=BoxBg,
  innerleftmargin=12pt,
  innerrightmargin=12pt,
  innertopmargin=8pt,
  innerbottommargin=8pt,
  skipabove=\baselineskip,
  skipbelow=\baselineskip,
]{customerquote}

\newmdenv[
  topline=false,
  bottomline=false,
  rightline=false,
  linewidth=3pt,
  linecolor=SectionBlue,
  backgroundcolor=BoxBg,
  innerleftmargin=12pt,
  innerrightmargin=12pt,
  innertopmargin=8pt,
  innerbottommargin=8pt,
  skipabove=\baselineskip,
  skipbelow=\baselineskip,
]{spokespersonquote}

\newcounter{faqnum} % lint-ok
\newenvironment{faqpair}[1]{%
  \needspace{4\baselineskip}%
  \vspace{0.5em}%
  \refstepcounter{faqnum}%
  \noindent\textbf{\color{SectionBlue}Q\thefaqnum: #1}\par\nopagebreak[4]%
  \vspace{0.2em}%
  \begin{adjustwidth}{1em}{0pt}%
  \setlength{\parindent}{0pt}%
  \setlength{\parskip}{0.5em}%
}{%
  \end{adjustwidth}%
  \vspace{0.5em}%
}
\newcommand{\faqref}[1]{\hyperref[#1]{FAQ~\ref*{#1}}}

\newcounter{featurenum}
\newcommand{\featureitem}[2]{%
  \refstepcounter{featurenum}%
  \item[\textbf{\color{SectionBlue}F\thefeaturenum.}]%
  \textbf{#1}: #2%
}
\newcommand{\featureref}[1]{\hyperref[#1]{Feature~\ref*{#1}}}

\usepackage{changepage}

% ══════════════════════════════════════════════════════════════════════════════
% Document
% ══════════════════════════════════════════════════════════════════════════════

\begin{document}

% ── Headline block ───────────────────────────────────────────────────────────
\begin{center} % lint-ok
  {\LARGE\bfseries\color{SectionBlue}
    vox Adds Podcasts and Audiobooks, Written With Your Coding Agent
    and Played While the Code Gets Written} \\[0.8em]
  {\large\color{AccentGray}
    Ask for a two-host series on Postgres internals, a noir mystery set in
    your billing service, or an hour about pandas, the animal. The agent
    writes it and vox voices it. The first part plays a couple of minutes
    later.} \\[0.3em]
  {\large\color{AccentGray}
    Free and open source from punt-labs. Music, shows and books share one
    library, one player and one session mood.}
\end{center}

\vspace{0.5em}

% ══════════════════════════════════════════════════════════════════════════════
% PRESS RELEASE
% ══════════════════════════════════════════════════════════════════════════════

\section*{Press Release}

SAN FRANCISCO, May 2027: Today, punt-labs announced podcasts and audiobooks
in vox, the free, open-source voice and audio layer for AI coding agents. A
developer working in Claude Code types \texttt{/podcast} or
\texttt{/audiobook} followed by whatever they want to hear. The agent in the
session writes the script and vox turns it into voiced audio. Playback
starts while the later parts are still being made. A podcast app plays
what somebody else recorded for somebody else
(see \faqref{faq:different}). Here the show is made inside the coding
session. The agent already knows the code and the day's work. The developer
steers it in plain conversation. Playback steps aside the moment the agent
needs attention (see \faqref{faq:interrupt}).

\prsection{Problem}

Coding with an agent is full of waiting. A typical turn takes under a
minute, but the longest one in a thousand runs past 45 minutes
\cite{anthropic2026agentautonomy}. Anthropic describes agent work that
spans hours \cite{anthropic2026longrunnagents}. The developer is present
for all of it and needed for little of it. Most people fill that time with
a phone, a chat thread or a feed. Each of those takes the eyes off the
terminal, and the agent's next question sits unanswered.

Audio is the obvious companion for half-attention, and vox already plays
generated music that follows the mood of the session. Music sets a mood. It
cannot teach you how the query planner works or tell you a story. For that,
developers open a podcast app, which knows nothing about the session. It
talks over the agent's permission prompt. It has no episode about the
library adopted this morning. It cannot be told ``less banter'' or ``go
deeper on that.'' Sometimes the developer wants the opposite of work:
twenty minutes about anything except the migration. The app answers with a
catalog to scroll.

\prsection{Solution}

\textbf{Ask for anything.} The developer types a wish in ordinary words:

{\small\begin{verbatim}
/podcast how are people testing agent loops? two hosts, one a skeptic
/audiobook tell me all about pandas. the animal, not the library
/audiobook a noir mystery set inside this codebase
/podcast a series called Learning React, start from zero
\end{verbatim}}

The agent answers with a short pitch in the conversation: a title, the cast,
and the list of episodes or chapters. Then it starts writing. There is no
form to fill in and no settings page. The topic can sit right next to the
work: the codebase, its domain, the technology under it. It can sit as far
from the work as the developer likes. Neither is the default.

\textbf{Listen while it is still being made.} Every show opens with a short
first part, about a minute long, so there is something to hear quickly. The
rest are generated ahead of the listener, a part or two in front. A long
book costs little more than what is heard (see \faqref{faq:cost}). A
podcast is a cast of hosts in conversation, the same voices every episode,
with a theme tune of its own. An audiobook is one narrator reading a book.
It runs as many chapters as the story needs.

\textbf{Steer it like a collaborator.} The show belongs to the developer and
the agent together. ``Less banter.'' ``Do the next episode on hooks.''
``That chapter dragged, redo it shorter.'' ``Tell the next one from the
villain's side.'' The agent rewrites what comes next. A series grows for
as long as there is something left to say (see \faqref{faq:steer}).

\textbf{It carries the vibe.} The mood of a session already drives the
agent's voice and the music. Shows and books take the same input. After a
morning of red test runs the hosts are gentler and the chapter is shorter.
On a good day they are quick and a little smug. The developer can also ask
for a mood outright, or for a language and a level: ``a B1 German audiobook
about a detective in Hamburg'' (see \featureref{feat:vibe} and
\featureref{feat:language}).

\textbf{It knows when to be quiet.} Two voices at once is the one thing
that does not work. When the agent has something to say, the show pauses
and the agent speaks. A few seconds after its last word, the show picks up
where it stopped. The developer can pause and resume at any time. Chimes
play over a show, as they play over music. Music never pauses. It already
plays at a lower level than the agent's voice. Every show and book keeps its place across sessions.
``Continue Learning React'' on Thursday resumes the sentence left on
Tuesday (see \faqref{faq:interrupt}).

\textbf{Everything is kept.} Shows and books are saved beside the music
albums. The player on the Lux display lists all three with their episodes
and chapters. Every script is saved as text next to its audio. Any of it
exports as tagged audio files for a walk or a commute
(see \faqref{faq:export}). The \texttt{vox}
command line plays, pauses, lists and exports the same library with no
agent in the loop.

\prsection{Customer Quote}

\begin{customerquote}
\textit{``I used to start a big refactor and then lose forty minutes to my
phone. Now I ask for something to listen to. Last week the agent and I made a
six-episode series on Postgres internals. The two hosts kept making jokes
about my actual schema. By the time the migration finished I was on episode
three and I finally understood what a visibility map is. On Friday the tests
had been red all afternoon, so I asked for a noir mystery set in our billing
service. The murderer was a cron job. I laughed out loud in an open office.
Then the agent chimed in to say the build was green, and the story carried
on from the same line.''}

\raggedleft \textbf{Priya Raman}, Staff Engineer at a logistics startup
\end{customerquote}

\prsection{Getting Started}

A developer who already has vox needs an ElevenLabs API key on a paid plan
and nothing else. First, update vox with the one-line installer. Second, in
any repository where vox is enabled, type \texttt{/podcast} or
\texttt{/audiobook} and say what to make. Third, keep working. The agent
posts its pitch within seconds and the first part plays about two minutes
later. There is no account to create with punt-labs, no upload step, and no
editor to learn.

\prsection{Spokesperson Quote}

\begin{spokespersonquote}
``We put the vibe in vibe coding. Literally. Music was the first half of
that. The agent plays DJ and the session has a mood you can hear. But I
kept wanting the other half. I sit next to something that can write about
anything, and much of my day is spent waiting for it. The best I could do
was open a podcast somebody else made for somebody else. So now I ask.
Sometimes it is a series about the thing I am building, because the agent
knows it better than any tutorial does. Sometimes it is pandas, the animal,
because I need ten minutes off. The design rule is that asking for a show
has to be as light as asking for a song. If it ever needs a form, we got it
wrong.''

\raggedleft \textbf{JM Freeman}, Creator of vox
\end{spokespersonquote}

\prsection{Call to Action}

Podcasts and audiobooks ship in vox, available now at
\url{https://github.com/punt-labs/vox}. The software is free and open
source under the MIT license. Generating audio uses the developer's own
ElevenLabs credits. Replaying anything already made is free. Update with
the one-line installer, restart Claude Code, and type \texttt{/podcast}.

\newpage

% ══════════════════════════════════════════════════════════════════════════════
% FREQUENTLY ASKED QUESTIONS
% ══════════════════════════════════════════════════════════════════════════════

\section*{Frequently Asked Questions}

\subsection*{External FAQs}

\begin{faqpair}{What is this and who is it for?}\label{faq:what}

  vox is the voice and audio layer for AI coding agents. It already gives
  the agent a voice, chimes, a session mood and generated background music.
  This release adds two things to make with the agent. A podcast is a cast
  of hosts talking through a subject over one or many episodes. An
  audiobook is a narrator reading a book of any length, chapter by chapter.

  It is for a developer who works beside a coding agent for hours and wants
  something good in their ears while the agent works. That can be something
  about the work or something that is a break from it.
\end{faqpair}

\begin{faqpair}{How is this different from NotebookLM's Audio Overviews, or from a podcast app?}\label{faq:different}

  NotebookLM turns a notebook of sources into a single audio overview. It
  can gather the sources itself from a topic prompt
  \cite{google2025deepresearch}. It offers four formats, a length setting and a steering prompt
  before generation. It has an English-only mode for asking the hosts
  questions \cite{notebooklm2026audiooverview,techcrunch2025notebooklmtone}.
  It is a good product and the closest thing to this one.

  \begin{center} % lint-ok
  \small
  \renewcommand{\arraystretch}{1.3}
  \begin{tabularx}{\linewidth}{@{} l X X @{}}
  \toprule
  \textbf{\color{SectionBlue}Dimension} & \textbf{\color{SectionBlue}NotebookLM} & \textbf{\color{SectionBlue}vox} \\
  \midrule
  Where it runs & A browser tab or phone app & The coding session \\
  What it knows & Sources you upload, or that it finds for a topic & The session: the code, the day's work, the mood \\
  Length & One overview per request & A series that grows, or a book with no length limit \\
  Subject & Non-fiction drawn from sources & Anything, including fiction \\
  Steering & A prompt before generation & Conversation before, during and after \\
  Interruptions & Unaware of your agent & Pauses for the agent and resumes \\
  Music & None & Generated albums in the same library \\
  \bottomrule
  \end{tabularx}
  \renewcommand{\arraystretch}{1.0}
  \end{center}

  A podcast app has better shows than any model will write this year. It
  loses on fit. It has no episode about your codebase, it cannot be
  redirected, and it keeps talking when the agent asks for approval.

  NotebookLM leads in two places. Its hosts have been in public use far
  longer. Its mode for asking the hosts questions is available today,
  and the vox equivalent is a follow-up (see \featureref{feat:call-in}).
\end{faqpair}

\begin{faqpair}{How do I get started?}\label{faq:getting-started}

  Update vox and check that an ElevenLabs key on a paid plan is configured.
  The \texttt{vox doctor} command reports this. Then type \texttt{/podcast}
  or \texttt{/audiobook} with a request. The agent replies with a title, a
  cast and a list of parts, then begins. The first part plays about two
  minutes later. The \texttt{pause}, \texttt{resume}, \texttt{next},
  \texttt{prev}, \texttt{list} and \texttt{status} subcommands work the way
  they do for \texttt{/music} today.
\end{faqpair}

\begin{faqpair}{How much does it cost?}\label{faq:cost}

  The software is free. Generating speech spends credits on your own
  ElevenLabs account. In a test call, 3,096 characters of text became three
  and a half minutes of narration \cite{voxspike2026audiobook}. At that
  rate an hour of narration is about 53,000 characters. The ElevenLabs list
  price for its current speech model is \$0.08 per 1,000 characters on
  every plan \cite{elevenlabs2026pricing}. An hour of narration therefore
  costs a little over four US dollars on any plan. That figure rests on one
  call. Your plan's balance is counted in credits, and the number of
  credits one character uses differs by plan, so the credit count you see
  will not match another account's.

  Two costs are unmeasured. Podcasts use the multi-speaker engine, which
  has not been priced by a test call. Writing the script spends model
  tokens on your coding agent's plan.

  Two rules keep the bill small. Audio is generated a part or two ahead of
  where you are listening. Abandoning a twelve-chapter book after chapter
  two costs three or four chapters, not twelve. Replaying anything already
  generated is free, forever. The \texttt{status} subcommand shows the
  credits each part used.
\end{faqpair}

\begin{faqpair}{What happens to my data?}\label{faq:data}

  The script is written by the agent in your session and sent to ElevenLabs
  to be voiced. If you asked for a show about your codebase, the script
  describes your codebase, and ElevenLabs receives that text. ElevenLabs
  keeps request data by default under its privacy policy. Its
  zero-retention mode is offered only to select enterprise customers
  \cite{elevenlabs2026zeroretention}. The script was written by your coding
  agent, so its model provider has seen it too. On your side, the audio and
  the script are saved on your own machine. There is no punt-labs cloud
  service. Usage telemetry is kept on your machine and sent nowhere unless
  you opt in. It records which commands, tools and hooks ran, which
  provider was used, whether each succeeded, whether the cache was hit and
  how long the text was. It never records the text itself.

  If text about your code must not leave your machine, do not ask for shows
  about your code. Shows about pandas are unaffected.
\end{faqpair}

\begin{faqpair}{Can it really be about anything? How accurate is it?}\label{faq:anything}

  Yes. Nothing ties a show to the repository. Technical requests work: a
  technology you are adopting, the domain your software serves, the
  architecture of the codebase in front of you. So do a language you are
  learning at a stated level, a story, and a subject with no connection to
  work at all.

  Accuracy is the accuracy of the agent that wrote it. A model writing from
  memory about a fast-moving subject will get things wrong in a confident
  voice. Two things help. The authoring instructions tell the agent to look
  things up before writing factual material, with whatever search and
  reading tools the session has. Every script is saved as text, so a claim
  that sounds off can be read, checked and corrected. Then that one part is
  regenerated.
\end{faqpair}

\begin{faqpair}{Won't people talking in my ears wreck my concentration?}\label{faq:focus}

  While you are reading or writing, yes, somewhat. Speech you are not
  trying to follow still interferes with keeping words in mind
  \cite{salame1982unattended}, and code is words. In one study, unattended
  speech hurt reading comprehension and instrumental music did not
  \cite{martin1988reading}. Music has a smaller cost and remains the
  default companion for heads-down work.

  Shows and books are for the other parts of the day: waiting on a long
  agent run, watching a build, skimming a diff, taking a break. One command
  pauses, and playback pauses by itself whenever the agent speaks.
\end{faqpair}

\begin{faqpair}{What happens when my agent needs me in the middle of an episode?}\label{faq:interrupt}

  The rule is that two voices never talk at once. If the agent speaks, the
  show pauses. It resumes from the same place a few seconds after the
  agent's last word. That delay is one setting. A run of several spoken
  updates becomes one pause, and you can lengthen or shorten the delay to
  taste. If the agent only chimes, the show keeps playing and the chime
  plays over it, as it does over music today. Music itself never pauses for
  the agent. It plays at a fixed lower level, so the agent's voice is heard
  over it.

  You stay in control. Pause and resume are one command each, and a pause
  you chose stays paused until you resume it. The player saves your
  position at least every 30 seconds while it plays. A crash or a restart
  costs you at most half a minute, and your place is kept across sessions.
\end{faqpair}

\begin{faqpair}{Can I change a show after it starts?}\label{faq:steer}

  Yes, by talking to the agent. You can ask for a different tone, more or
  less depth, a new episode on a named subject, a chapter retold, or a host
  replaced. Parts not yet generated are rewritten before they are voiced.
  A part already voiced can be regenerated by itself without touching the
  rest. A series has no fixed length. ``Add an episode on hooks'' three
  weeks later appends one, and the hosts remember what they already
  covered.
\end{faqpair}

\begin{faqpair}{Can I take a show with me?}\label{faq:export}

  Yes. Any show, book or album exports as ordinary audio files with title,
  part number and series tags, ready for a phone or a car. The scripts
  export with them.
\end{faqpair}

% ── Internal FAQs ────────────────────────────────────────────────────────────

\newpage
\subsection*{Internal FAQs}

\subsubsection*{Value \& Market}

\begin{faqpair}{What is the total addressable market?}\label{faq:tam}

  Small, and not the reason to build this. The audience is a subset of the
  audience in the vox product PR/FAQ. That document estimates 360,000 to
  720,000 developers who run agent tasks long enough to warrant a
  notification \cite{voxprfaq2026}. Its starting count of developers is an
  uncited assumption. This feature narrows that further to people with a
  paid ElevenLabs plan who want long-form listening at their desk. No
  source sizes that group.

  The reason to build it is that it makes vox the audio side of working
  with an agent (see \faqref{faq:different}). Adoption will follow from
  that or it will not. The first market is the people who build vox.
\end{faqpair}

\begin{faqpair}{What evidence do we have that developers want this?}\label{faq:evidence}

  None from customers. The evidence is indirect.

  \textbf{The category exists.} Google keeps adding to its generated
  podcasts: new formats in 2025, then source gathering from a topic prompt
  \cite{techcrunch2025notebooklmtone,google2025deepresearch}. The inference
  is that people listen to machine-written two-host audio when it is about
  something they chose. No usage data backs that inference.

  \textbf{The slot may exist.} The longest agent turns run past 45 minutes,
  but those are one turn in a thousand \cite{anthropic2026agentautonomy}.
  The share of a session spent waiting is unmeasured. Claude Code
  transcripts on the builders' machines have timestamps and can supply it.

  \textbf{One user.} vox music is in daily use by its author, who asked for
  this. One person's wish is an anecdote.

  \textbf{What is missing.} Nobody has listened to a vox-made episode,
  because none exists. Whether an agent-written show keeps a listener past
  the first novelty is unknown. That is the main value risk
  (see \faqref{faq:failure}).
\end{faqpair}

\begin{faqpair}{Who are the competitors and why would we win?}\label{faq:competitors}

  \textbf{NotebookLM.} The closest product, covered in
  \faqref{faq:different}. Google built a standalone API that turns supplied
  content into a podcast file with no notebook. It is now deprecated and
  closed to new customers \cite{google2026podcastapi}. If Google reopened
  it, a coding agent could call it and get polished two-host audio about a
  set of files. It would still lack the session. That means playback that
  yields to the agent, a library shared with music, a mood, a series that
  grows by conversation, and books.

  \textbf{ElevenLabs.} It owns the engines vox calls and sells its own
  long-form tools. Its Studio API can create a two-voice podcast from
  source material \cite{elevenlabs2026studiopodcast}. The Studio API refused this account
  with a 403 \cite{voxspike2026audiobook}. Its reader app generates
  podcasts from a user's documents \cite{elevenlabs2024genfm}. ElevenLabs
  could offer agent-made shows directly, and it would do the audio better
  than vox. It would still lack the session.

  \textbf{Podcast and audiobook apps.} Better content, no fit. They are the
  incumbent for the listener's time. They win by default if generated shows
  are dull.

  \textbf{Other agent voice tools.} AgentVibes plays a bundled music
  catalog \cite{voxprfaq2026}.

  The position is narrow: vox wins only on being in the session. If that
  turns out not to matter to listeners, nothing else here is defensible.
\end{faqpair}

\begin{faqpair}{How will people find it?}\label{faq:acquisition}

  Existing vox users get both commands with an update. Beyond them, the
  channel is a demo. A sixty-second clip of a mystery set in a well-known
  open-source codebase is easy to share and explains the feature by
  itself. There is no paid acquisition.
\end{faqpair}

\begin{faqpair}{What is your next step to validate the vision?}\label{faq:validation-plan}

  Listen. The API probe produced a three-and-a-half-minute narrated chapter
  and a two-voice dialogue clip, and nobody has played either
  \cite{voxspike2026audiobook}. The next step is one sitting. Play both
  clips. Have the agent write one ten-minute podcast episode and one
  five-minute chapter. Generate them through the probe script, extended to
  take a script file. Listen to them during an ordinary agent run. Invite
  three to five developers from outside punt-labs to listen too.

  That sitting answers three questions about the content. Do the hosts
  sound like two people or one model taking turns? Can you hear the joins
  where dialogue requests are stitched? Is a ten-minute agent-written
  episode worth ten minutes? If the answer to the last one is no, the work
  stops before any product code is written.

  The same sitting takes seven measurements.

  \begin{itemize}[leftmargin=1.5em, nosep]
    \item How long the agent takes to write a five-minute chapter.

    \item How long one chapter takes to generate, streamed and not
          streamed.

    \item Whether one 20,000-character speech request is accepted.

    \item Whether the multi-speaker engine accepts a request linked to the
          one before it.

    \item Whether a long coding task slows or stalls while a background
          agent writes a script beside it.

    \item Whether the background agent's output returns into the coding
          session's context.

    \item Whether a request typed during a long agent turn is answered at
          once or queued until the turn ends.
  \end{itemize}
\end{faqpair}

\subsubsection*{Technical}

\begin{faqpair}{What are the major technical risks?}\label{faq:tech-risks}

  \textbf{Who writes the next part during a run (high).} The agent that
  writes the script is the agent doing the coding. A session takes one turn
  at a time, so a script written in the main session interrupts the coding
  or waits for it. The script would also fill the coding agent's context.
  The design answer is a background agent with its own context. The main
  agent starts it when the show is requested. It writes the outline and the
  parts and hands each finished script to the daemon. The daemon has no
  model and cannot ask for part four, so the writer has to stay ahead of
  the listener unprompted. It learns the listener's position by asking the
  daemon, and each check costs tokens. The writer ends when the session
  ends. The show then plays what was already written and stops. A later
  session starts a new writer from the outline and running summary stored
  with the show (see \featureref{feat:memory}). Today a chime plays when
  any sub-agent starts or stops, and the writer must be exempt. None of
  this has been tried. Four things are unknown. Does the coding task slow?
  Does the writer's output return into the coding session's context? How
  many model tokens does the writing spend? What does the listener hear
  when the writer falls behind?

  \textbf{Stitched dialogue (high).} ElevenLabs recommends at most 2,000
  characters of text per multi-speaker request
  \cite{elevenlabs2026dialogue}. The measured narration rate is about 880
  characters per minute \cite{voxspike2026audiobook}. Dialogue pacing may
  differ, so applying that rate to a podcast is an assumption. On that
  assumption a ten-minute episode is about 8,800 characters, or five
  requests joined end to end. Linking one request to the next was accepted
  for narration. It has not been tried on the multi-speaker engine, and no
  join has been heard. NotebookLM, with far more tuning, still warns of voice
  glitches and speakers swapping \cite{notebooklm2026audiooverview}.
  Mitigation: listen first (see \faqref{faq:validation-plan}). If joins are
  audible, the fallback is one voiced turn per request through the
  single-speaker engine. At roughly 90 characters a turn, that is about a
  hundred requests per episode, with a risk of tone breaks between turns.

  \textbf{Chapter size (medium).} Narration is the easy engine: one request
  per chapter. ElevenLabs' model documentation lists a 100,000-character
  limit per request for the current model \cite{elevenlabs2026models}. Its
  pricing page says 10,000 \cite{elevenlabs2026pricing}. The probe sent
  3,096. A five-minute chapter is about 4,400 characters and fits under
  either figure. Longer chapters are unproven until a larger request is
  sent.

  \textbf{Writing a long book (medium).} A ten-hour book is roughly half a
  million characters, far more than an agent writes in one turn. The open
  design question is when chapter seven gets written. Writing all the text
  at creation makes creation a long job. Writing each chapter as the
  listener approaches needs an agent available days later. The design
  assumes the outline and a running summary are stored with the book. Any
  later session can then write the next chapters from them. That
  assumption is untested.

  \textbf{Sharing the speakers (medium).} Today music plays at a fixed lower
  level than speech and chimes, and that works well. Music keeps that
  behaviour. Chimes play over a spoken show too. Agent speech over a spoken
  show is the one case that fails. The show pauses for agent speech and for
  a voice call, then resumes after a set delay. That delay is the one value
  to tune, and it is a user setting. During playback the daemon reports
  its position as a heartbeat. The saved position and the progress bar both
  come from that heartbeat. A saved position in a local store is routine
  work. The daemon's playback state gains six things. Three are pauses: one
  for agent speech, one the developer chose, and a wait for a part that is
  not ready. The others are a timed resume, a saved position per show, and
  an end state. That is a state machine change and gets a formal model before
  code.

  \textbf{Time to first audio (unknown).} The target is audio within about
  two minutes of the request. Nothing has been timed. The wait is the
  agent's writing time plus generation time. The design shortens both by
  opening every show with a part about a minute long. The listening session
  times each step. If the target cannot be met, the opening part gets
  shorter or music from the library plays during the wait.

  \textbf{One vendor (medium).} Every engine here is ElevenLabs. Its
  production API for long-form work is closed to this account and answers
  403. That is why vox calls the plain speech and dialogue endpoints
  directly \cite{voxspike2026audiobook}. A pricing or model change lands on
  this feature immediately. Single-narrator audiobooks could run on OpenAI
  or Polly as a hedge (see \featureref{feat:other-providers}).
\end{faqpair}

\begin{faqpair}{What does this depend on?}\label{faq:dependencies}

  \textbf{Built already.} The library model that music uses: a named,
  saved collection of ordered parts, each with its own file. Around it sit
  a catalog, a filler that generates ahead of playback, tagged export, the
  Lux player, and the command-line and agent tool paths. Podcasts and
  audiobooks are two more kinds of the same thing.

  \textbf{ElevenLabs.} Long-form speech, multi-speaker dialogue, and music
  for theme tunes, on a paid plan.

  \textbf{The agent.} All writing is done by the model in the session. vox
  ships instructions for each format and has no model of its own. Output
  quality tracks the model the developer runs.

  \textbf{Not needed.} No punt-labs service, no account, no upload.
\end{faqpair}

\begin{faqpair}{What is the delivery order, and what gets cut?}\label{faq:timeline}

  In order:

  \begin{enumerate}[leftmargin=1.5em, nosep]
    \item \textbf{Listening session.} This decides go or stop
          (see \faqref{faq:validation-plan}).

    \item \textbf{Playback for spoken shows.} This covers in-order play, an
          end, pausing for agent speech, resuming after a set delay, a
          position heartbeat, a saved position and a progress bar. Music
          gets the progress bar too. It is modeled formally first and is
          testable with any existing audio.

    \item \textbf{Background writer and show memory.} Scripts are written
          off the main turn, from an outline and running summary stored
          with the show. Done means a long coding task and a show run side
          by side, and a new session can write the next part.

    \item \textbf{Audiobook, one narrator.} The simplest engine path, with
          a one-minute opening part and each script saved as text. The
          request can set a mood, a language and a level. Done means: ask,
          hear the opening, keep listening as later chapters arrive, stop
          and resume next day.

    \item \textbf{Podcast, two hosts.} The dialogue engine, a stable cast
          and a theme tune, with the same mood and language inputs. Done
          means the same test as the audiobook, with a series of three
          episodes.

    \item \textbf{Steering.} Append an episode, regenerate one part, and
          rewrite what has not been voiced yet.

    \item \textbf{Library and export for both formats.} The Lux player,
          per-part status, the command line and tagged export. Most of this
          exists for music.
  \end{enumerate}

  Items 2 to 7 cover the Must Do list and ship as one release. If scope
  compresses, every Should Do item waits, starting with character voices.
  If the release has to be split, items 2 to 4 go first as audiobooks, and
  podcasts follow. No dates are given because nothing comparable has been
  built here yet to estimate from.
\end{faqpair}

\begin{faqpair}{What breaks first if people use this a lot?}\label{faq:scaling}

  The listener's bill, then the disk. An hour of narration every working
  day is about 1.2 million characters a month. At list price that is a
  little over ninety dollars \cite{elevenlabs2026pricing}. The entry plans
  include far fewer characters than that, so a daily listener needs a
  larger plan or pays for overage. A hobbyist will notice. Audio files
  accumulate at about a megabyte a minute. Neither is a vox scaling
  problem. Both need to be visible: credits per part in \texttt{status},
  and library size in \texttt{vox doctor}.
\end{faqpair}

\subsubsection*{Business}

\begin{faqpair}{What is the revenue model?}\label{faq:revenue}

  There is none. vox is free under the MIT license and this feature
  changes nothing about that. The strategic value is identity. A tool that
  makes a show with you is a different thing from a tool that reads
  notifications aloud. It is the kind of thing people show each other.
\end{faqpair}

\begin{faqpair}{What does it cost us?}\label{faq:pl}

  Maintainer time and test credits. Running costs are zero because vox has
  no service. Every generated minute is paid for by the developer who asked
  for it. Live testing of long-form audio is the new expense. A full pass
  over both formats spends a few dollars of credits. A notification test
  spends a fraction of a cent.
\end{faqpair}

\begin{faqpair}{How will we know it worked?}\label{faq:metrics}

  Usage telemetry stays on each listener's machine unless they opt in, so
  the measures are things a person can report.

  \begin{itemize}[leftmargin=1.5em, nosep]
    \item \textbf{We use it.} After two weeks, the people building vox
          still choose a vox-made show over their podcast app during agent
          runs, unprompted.

    \item \textbf{Finished parts.} At least 70\% of parts that start
          playing reach their end. The daemon records this in the library
          on the listener's own machine. Each listener reads the figure
          from \texttt{status} and reports it.

    \item \textbf{Second requests.} Someone who makes one show makes
          another within a week.

    \item \textbf{The demo test.} A developer shown a mystery set in their
          own codebase asks how to install it.
  \end{itemize}

  The decision threshold is two weeks of daily availability. If the
  builders have gone back to their podcast apps by then, the content is not
  good enough. Work stops at whatever has shipped.
\end{faqpair}

\begin{faqpair}{Why now?}\label{faq:why-now}

  Three things are now true. One speech request on the current model
  produced three and a half minutes of narration
  \cite{voxspike2026audiobook}. The documented limit per request is at
  least 10,000 characters \cite{elevenlabs2026models,elevenlabs2026pricing}. The
  multi-speaker engine returns audio on the same model
  \cite{voxspike2026audiobook}. And the music playlist in vox became a
  general library of saved, ordered, replayable parts. Two new formats are
  additions to something that exists.
\end{faqpair}

\begin{faqpair}{What are we not building?}\label{faq:not-building}

  A publishing tool, an audio editor, or a notebook product. There is no
  hosting or distribution of shows. There is no timeline for trimming
  audio. Nobody is asked to upload sources. There is no language model
  inside vox and no imitation of an actual person's voice.
\end{faqpair}

\begin{faqpair}{What does failure look like?}\label{faq:failure}

  \textbf{It is dull.} The likeliest failure. The first show is a delight
  and the fourth sounds like the first: the same two cheerful voices, the
  same structure, the same ``great question.'' Listeners go back to shows
  made by people. Response: the per-format authoring instructions are the
  product. They get the iteration that the audio pipeline does not need.
  They are where vibe, cast and form turn into variety.

  \textbf{It is wrong.} A confident host states something false about a
  subject the listener is trying to learn. One such moment costs the
  format its trust for teaching, leaving only fiction. Response: lookup
  before writing, scripts kept as text, and a plain statement in the
  documentation that this is model-written material.

  \textbf{It is in the way.} Speech fights with reading
  (see \faqref{faq:focus}) and pausing is fiddly. People conclude that
  music was the right amount of audio. Response: this is acceptable.
  Music stays, and the two formats remain available for the waiting.
\end{faqpair}

% ══════════════════════════════════════════════════════════════════════════════
% FOUR RISKS ASSESSMENT
% ══════════════════════════════════════════════════════════════════════════════

\newpage
\section*{Risk Assessment}

\renewcommand{\arraystretch}{1.6} % lint-ok
\begin{tabularx}{\textwidth}{@{} l l X @{}}
\toprule
\textbf{\color{SectionBlue}Risk} & \textbf{\color{SectionBlue}Rating} & \textbf{\color{SectionBlue}Assessment} \\
\midrule
Value      & \textcolor{RiskRed}{\textbf{High}} &
  No one has heard a vox-made show. The bet is that a show written for you,
  about what you chose, inside your session, beats a better-produced show
  that is none of those things. Novelty will cover the first three
  episodes. Nothing yet says what covers the tenth. \\
Usability  & \textcolor{RiskAmber}{\textbf{Medium}} &
  Asking is one line and matches \texttt{/music}, which people already use.
  The risk is in the listening. Speech competes with reading, so the
  feature fits only part of the day. Agent turns are often under a minute,
  so a talkative agent pauses the show often. The resume delay turns a
  burst of updates into one pause. The listener can change that delay,
  pause by hand, or switch the agent to chimes. Pause, resume and
  place-keeping have to work every time. \\
Feasibility & \textcolor{RiskRed}{\textbf{High}} &
  The library, the player and the generation pipeline exist for music. The
  API calls were made and returned audio. Four things are unproven, and the
  first has no fallback. Can a
  background agent write scripts beside a coding task? How does stitched
  dialogue sound? How long does generation take? How does a book longer
  than one session get written? \\
Viability  & \textcolor{RiskAmber}{\textbf{Medium}} &
  No revenue to validate and no service to run. Two exposures remain. Every
  engine comes from one vendor. A daily listener pays about ninety dollars
  a month for something that competes with free podcasts. \\
\bottomrule
\end{tabularx}
\renewcommand{\arraystretch}{1.0}

% ══════════════════════════════════════════════════════════════════════════════
% FEATURE APPENDIX
% ══════════════════════════════════════════════════════════════════════════════

\newpage
\section*{Feature Appendix}

\subsection*{Must Do}

\begin{enumerate}[nosep,leftmargin=2.5em]
  \featureitem{\texttt{/podcast} and \texttt{/audiobook}}{Two commands, each with its own instructions to the agent on how to write that form. The request is free text. The agent replies with a pitch (title, cast, parts) and starts.}\label{feat:commands}

  \featureitem{Open subject}{Nothing ties a show to the repository. Code, domain, technology, fiction, or a subject unrelated to work are all ordinary requests.}\label{feat:open-subject}

  \featureitem{Vibe as an input}{The session mood that already drives voice and music is passed to the agent when it writes. It affects tone, pace and length.}\label{feat:vibe}

  \featureitem{Language and level}{A show can be requested in a target language at a stated proficiency level.}\label{feat:language}

  \featureitem{Podcast form}{A stable, named cast of hosts in conversation. Episodes stand alone and open with a short recap. A series of one is normal.}\label{feat:podcast}

  \featureitem{Theme tunes}{Each podcast gets one short generated intro and outro from the music engine, made once per series and mixed into every episode.}\label{feat:theme}

  \featureitem{Audiobook form}{One narrator, chapters that follow one another, no limit on length.}\label{feat:audiobook}

  \featureitem{Background writer}{Scripts are written by a background agent with its own context. The coding session's turn and context stay free.}\label{feat:writer}

  \featureitem{Play while generating}{Every show opens with a part about a minute long, which plays when ready. Later parts are generated a part or two ahead of the listener and no further.}\label{feat:lazy}

  \featureitem{In-order playback with an end}{Spoken shows play in sequence and stop after the last part. Music keeps rotating.}\label{feat:sequential}

  \featureitem{Show memory}{Each show stores its outline and a running summary of what has been covered. A later session starts a writer from them, so a series continues days later and the hosts do not repeat themselves.}\label{feat:memory}

  \featureitem{Yield to speech}{A spoken show pauses for agent speech and voice calls. It resumes from the same position a set time after the last spoken word. That delay is a user setting. Chimes play over the show. Music never pauses. A pause the developer chose never resumes by itself.}\label{feat:yield}

  \featureitem{Playback heartbeat and progress bar}{While anything plays, the daemon reports the current part and position at least every 30 seconds. The Lux player draws a progress bar from it for music, shows and books.}\label{feat:progress}

  \featureitem{Keep my place}{The heartbeat position is saved in the local library for every show and book. It survives sessions and daemon restarts. ``Continue'' resumes it.}\label{feat:bookmark}

  \featureitem{One library}{Shows, books and albums are saved in one catalog. One player on the Lux display lists them with titles, parts and progress. The \texttt{vox doctor} command reports the library's size on disk.}\label{feat:library}

  \featureitem{Per-part status}{Each part reports pending, generating, ready or failed with a reason, and the credits it used.}\label{feat:status}

  \featureitem{Scripts kept as text}{The written script of every part is stored beside its audio and can be read.}\label{feat:scripts}

  \featureitem{Steering}{Rewrite parts not yet voiced. Append a part to an existing show. Regenerate one voiced part without touching the others.}\label{feat:steering}

  \featureitem{Command-line parity}{\texttt{vox podcast} and \texttt{vox audiobook} play, pause, list, show status and export with no agent. They also accept a finished script on standard input and generate from it.}\label{feat:cli}

  \featureitem{Export}{Any show or book exports as tagged audio files with its scripts.}\label{feat:export}
\end{enumerate}

\subsection*{Should Do}

\begin{enumerate}[nosep,leftmargin=2.5em]
  \featureitem{Character voices}{Audiobook dialogue lines spoken by a cast of distinct voices through the multi-speaker engine.}\label{feat:characters}

  \featureitem{Call in to the show}{Pause an episode and ask the hosts a question by voice, using the existing vox call. They answer in character and the episode resumes.}\label{feat:call-in}

  \featureitem{The agent offers}{When the session mood has been low for a while, the agent suggests a break and a show to go with it. Off by default.}\label{feat:offers}

  \featureitem{Designed voices}{A host or character voice created from a description (``older, dry, Glaswegian'') and kept with the show.}\label{feat:designed-voices}

  \featureitem{Score and sound effects}{A music bed and effects mixed into audiobook chapters at generation time.}\label{feat:score}

  \featureitem{Other speech providers}{Single-narrator audiobooks through OpenAI or Polly for listeners without ElevenLabs.}\label{feat:other-providers}

  \featureitem{Progress in the terminal}{The same heartbeat drawn as a progress line in a terminal view, for sessions with no Lux display.}\label{feat:progress-tui}

  \featureitem{Cost preview}{Before a long work starts, the pitch states an estimated credit cost per part.}\label{feat:cost-preview}
\end{enumerate}

\subsection*{Won't Do}

\begin{enumerate}[nosep,leftmargin=2.5em]
  \featureitem{Publishing}{No feeds, no hosting, no upload to podcast directories. Shows are made for the person who asked. Export covers taking a file elsewhere.}\label{feat:no-publishing}

  \featureitem{An audio editor}{No timeline, trimming or per-sentence regeneration. The unit of change is a part, and the way to change it is to ask.}\label{feat:no-editor}

  \featureitem{Source upload}{No notebook of documents to manage. The agent reads what the session can already reach.}\label{feat:no-upload}

  \featureitem{Two programs at once}{No spoken show over a separately playing music program. Any music under speech is mixed into the part when it is generated.}\label{feat:no-mixing}

  \featureitem{Cloning voices}{No hosts that imitate the voice of an actual person. That needs the person's consent, and vox has no way to check it.}\label{feat:no-cloning}

  \featureitem{A built-in writer}{There is no language model inside vox. The agent in the session writes, and the command line takes finished scripts.}\label{feat:no-llm}
\end{enumerate}

% ══════════════════════════════════════════════════════════════════════════════
% REFERENCES
% ══════════════════════════════════════════════════════════════════════════════

\newpage
{\emergencystretch=1em
\printbibliography[title={References}]
}

\end{document}
